\chapter{The Alchemist's Workbench: Practical Debugging, REPL Mastery, and Emacs Workflows}
\label{chap:debugging-and-repl}

\epigraph{``If you cannot inspect the living tissue of your operating system while it breathes, you are not administering a computer; you are merely performing archaeology on its corpses.''}{--- The Lisp Sysadmin}

\noindent
Up to this point in our journey, we have celebrated the declarative serenity of GNU Guix. We have declared operating systems in single files, constructed hermetic development shells, traveled through time with channel pins, and packaged software with functional purity.

And yet, every developer eventually encounters that fateful afternoon: you add a new service to your \texttt{/etc/config.scm}, or write a custom package definition for an internal library, and type:
\begin{lstlisting}[style=guixcli]
sudo guix system reconfigure /etc/config.scm
\end{lstlisting}

Instead of instantaneous enlightenment, your terminal erupts into a 40-line cascade of Guile Scheme stack traces:
\begin{lstlisting}[style=guixcli]
Backtrace:
In ice-9/boot-9.scm:
  1752:10 18 (with-exception-handler _ _ #:unwind? _ # _)
In gnu/services.scm:
   884:32 17 (loop _ #f)
In ice-9/boot-9.scm:
  1685:16 16 (raise-exception _ #:continuable? _)
  1685:16 15 (raise-exception _ #:continuable? _)
ice-9/boot-9.scm:1685:16: In procedure struct-vtable: Wrong type argument: #f
\end{lstlisting}

\begin{tearsbox}[The Imperative Developer in the Land of Parentheses]
If you come from an imperative background, your standard debugging instincts immediately trigger:
\begin{itemize}
  \item In Python or Node.js, you would scatter \texttt{print("HERE 1")} and \texttt{console.log()} statements across the codebase.
  \item In Docker, you would run \texttt{docker run -it broken-image /bin/bash} to poke around the filesystem interactively.
  \item In Ubuntu or Arch, you would run \texttt{strace -f}, inspect \texttt{/var/log/syslog}, or rebuild packages with manual \texttt{./configure \&\& make} invocations.
\end{itemize}
In Guix, none of these habits seem to work at first glance! The build runs inside an isolated chroot container that vanishes when a failure occurs. The configuration is evaluated by a Scheme daemon, not executed as a Bash script. And that terrifying backtrace references \texttt{ice-9/boot-9.scm}---a core Guile internal file you didn't write and certainly don't want to debug on a Friday afternoon!
\end{tearsbox}

Here is the liberating truth: \textbf{GNU Guix is not a static tool parsing dead configuration files. It is an interactive, introspectable Lisp environment.} 

Because Guix is built on GNU Guile Scheme, you have access to interactive REPLs, structural record reflection, step-by-step debuggers, instant source navigation, and live system inspection that leave traditional Unix distributions decades behind. 

In this chapter, we transform you from a confused observer deciphering cryptic backtraces into a \textbf{Diagnostic Alchemist} equipped with the complete Guix debugging workbench.

\section{The Three Diagnostic Domains: Where Did It Fail?}

The first rule of Guix debugging is identifying \textbf{which stage} of the evaluation lifecycle failed. Errors in Guix fall into three distinct domains:

\begin{enumerate}
  \item \textbf{Evaluation Time (Scheme Evaluation)}:
    \begin{itemize}
      \item \textit{When}: Happens while Guix is reading and evaluating your Scheme code (\texttt{config.scm}, \texttt{manifest.scm}, or package recipe) to compute the dependency graph.
      \item \textit{Symptoms}: Guile backtraces, syntax errors, unbound variables, wrong record types, or cycle errors in the service extension graph.
      \item \textit{Tooling}: \guixcmd{guix repl}, Guile metacommands, Geiser, and record accessors.
    \end{itemize}

  \item \textbf{Build Time (Derivation Execution inside Daemon Sandbox)}:
    \begin{itemize}
      \item \textit{When}: Scheme evaluation succeeded and produced a derivation (\texttt{.drv}), but compiling or installing the software inside the isolated \texttt{/tmp/guix-build-...} container failed.
      \item \textit{Symptoms}: Compiler errors (\texttt{gcc: error: ld returned 1 exit status}), failed unit tests, missing patch targets, or network errors during tests.
      \item \textit{Tooling}: \guixcmd{guix build -K} (\texttt{--keep-failed}), \texttt{--log-file}, and interactive chroot reproduction via \guixcmd{guix shell -D}.
    \end{itemize}

  \item \textbf{Runtime (Shepherd Init and Daemons)}:
    \begin{itemize}
      \item \textit{When}: The system built and reconfigured successfully, but a system service (e.g. Nginx, PostgreSQL, OpenSSH) crashes or refuses to start at boot.
      \item \textit{Symptoms}: \texttt{herd status} shows a service in \texttt{stopped} state, or network daemons fail to bind ports.
      \item \textit{Tooling}: \guixcmd{herd detailed-status}, \guixcmd{herd doc}, \texttt{/var/log/shepherd.log}, and ephemeral virtual machines (\guixcmd{guix system vm}).
    \end{itemize}
\end{enumerate}

\section{CLI Triage: Crime Scene Investigation in the Store}

When a build fails, the \texttt{guix-daemon} strictly adheres to the principle of cleanliness: it deletes the temporary build sandbox to prevent disk clutter. This is wonderful for disk hygiene, but infuriating when you need to know why \texttt{make check} threw a core dump.

\subsection{Saving the Crime Scene: \texttt{guix build -K}}

The single most valuable command-line flag for package authors and maintainers is \guixcmd{-K} (or \guixcmd{--keep-failed}):

\begin{lstlisting}[style=guixcli]
# Keep the build directory intact upon error
guix build -K my-broken-package
\end{lstlisting}

When the build fails, Guix prints the exact path of the preserved build directory:
\begin{lstlisting}[style=guixcli]
builder for `/gnu/store/...-my-broken-package-1.0.drv' failed with exit code 1
note: keeping build directory `/tmp/guix-build-my-broken-package-1.0.drv-0'
\end{lstlisting}

\begin{protipbox}[Anatomy of a Preserved Build Chroot]
Inside \texttt{/tmp/guix-build-<name>.drv-0}, you will find the exact state of the universe at the moment of failure:
\begin{itemize}
  \item \texttt{environment-variables}: A flat shell script containing every single environment variable (\texttt{PATH}, \texttt{CPATH}, \texttt{LIBRARY\_PATH}, \texttt{PKG\_CONFIG\_PATH}) set by the build system.
  \item The extracted upstream source code, intermediate object files (\texttt{.o}), and build artifacts.
  \item \texttt{config.log} or \texttt{CMakeError.log}: Detailed compiler probes showing why configure scripts failed.
  \item \texttt{test-suite.log}: The complete stdout/stderr of failed unit test suites.
\end{itemize}
\end{protipbox}

\subsection{Re-entering the Build Environment Interactively}

Don't just inspect the files---re-enter the build environment and run the commands yourself! Using \guixcmd{guix shell} with the \texttt{-D} (development environment) flag, you can spawn an interactive shell with every dependency present, source the environment variables, and run a debugger:

\begin{lstlisting}[style=guixcli]
# 1. Step into an interactive shell with all build inputs exposed
guix shell -D my-broken-package -- bash

# 2. Navigate to the preserved build directory
cd /tmp/guix-build-my-broken-package-1.0.drv-0/source

# 3. Source the exact build environment captured by the daemon
source ../environment-variables

# 4. Run the failing step directly with gdb or verbose flags
make test
# or
gdb --args ./tests/failing-test-binary
\end{lstlisting}

You have now bridged the gap between pure functional containers and the raw interactive power of GNU/Linux. Once you fix the bug, exit the shell and delete the directory in \texttt{/tmp}.

\subsection{Inspecting Build Logs: \texttt{guix build --log-file}}

Never squint at a fast-scrolling terminal to catch a compiler warning. Every build log is compressed and preserved in \texttt{/var/log/guix/drvs/}:

\begin{lstlisting}[style=guixcli]
# View the build log of any package by name
guix build --log-file nginx | less

# Or view the log of a specific derivation path
guix build --log-file /gnu/store/5q3...-nginx-1.24.0.drv | grep -i error
\end{lstlisting}

If a build was downloaded as a pre-built binary substitute from a continuous integration server (such as \texttt{ci.guix.gnu.org}), Guix will fetch the remote build log over the network automatically!

\subsection{Forensic Dependency Graphing: \texttt{guix graph} and \texttt{guix size}}

Have you ever wondered: \emph{``Why is my tiny 2\,MB CLI tool pulling in 1.4\,GB of dependencies including Qt5 and Python?''}

Before tearing your hair out, use \guixcmd{guix size} and \guixcmd{guix graph} to locate the culprit:

\begin{lstlisting}[style=guixcli]
# Inspect the closure size breakdown of a package
guix size my-tool
\end{lstlisting}

Guix prints a beautiful table showing every package in the closure, its self-size, its total dependency closure, and the percentage of the overall disk space it occupies:
\begin{lstlisting}[style=guixcli]
store item                               total-size (MiB)  self-size (MiB)
/gnu/store/...-gcc-11.3.0-lib                 74.2              74.2   45.2%
/gnu/store/...-glibc-2.35                     40.1              40.1   24.4%
/gnu/store/...-my-tool-1.0                    164.1              2.1    1.3%
\end{lstlisting}

To find the exact link connecting your package to an unwanted dependency:
\begin{lstlisting}[style=guixcli]
# Output an SVG dependency graph showing the shortest path to an input
guix graph --path my-tool unwanted-dependency | dot -Tsvg > path.svg
\end{lstlisting}
You will immediately discover whether the dependency was introduced via an accidental \texttt{propagated-inputs} declaration or a hidden runtime path baked into an executable's \texttt{RUNPATH}.

\section{Source Navigation: Jumping Directly to Definitions}

One of the greatest joys of GNU Guix is that its package recipes and system services are not hidden behind proprietary databases, obfuscated binary manifests, or disparate shell scripts. They are readable Scheme files sitting directly on your filesystem.

\subsection{The Instant CLI Miracle: \texttt{guix edit}}

Want to see how the official Guix maintainers packaged \texttt{postgresql}, \texttt{emacs}, or \texttt{linux-libre}? You don't need to open a web browser or search GitHub:

\begin{lstlisting}[style=guixcli]
guix edit postgresql
\end{lstlisting}

\guixcmd{guix edit} queries Guile's internal source properties, resolves the file and line number where \texttt{postgresql} was defined via \texttt{define-public}, and immediately opens your configured \texttt{\$EDITOR} (or \texttt{\$VISUAL}) at that exact line!

\begin{lstlisting}[style=guixcli]
# Jump to the definition of a package in your favorite editor
EDITOR="emacsclient -c" guix edit nginx
# or
EDITOR="nvim" guix edit rust
\end{lstlisting}

\subsection{Locating Services and System Modules}

While \guixcmd{guix edit} works for package recipes, Guix System services live inside the \texttt{gnu/services/} module tree. You can find the exact filesystem location of any Guix module using Scheme's \texttt{search-path} inside a one-line evaluation:

\begin{lstlisting}[style=guixcli]
# Find the exact path of the desktop services module
guix repl -e '(display (search-path %load-path "gnu/services/desktop.scm"))'
\end{lstlisting}

\begin{alchemybox}[The Guix Source Code is the Ultimate Manual]
The official Guix documentation is extensive, but the source code itself is often even clearer:
\begin{itemize}
  \item Package definitions live in \texttt{gnu/packages/<category>.scm} (e.g. \texttt{gnu/packages/web.scm}, \texttt{gnu/packages/python.scm}).
  \item System services live in \texttt{gnu/services/<category>.scm} (e.g. \texttt{gnu/services/ssh.scm}, \texttt{gnu/services/networking.scm}).
  \item Service configuration records are defined using \texttt{define-configuration}. Looking at the source file shows you every configurable field, its default value, and its accepted types in 10 seconds flat!
\end{itemize}
\end{alchemybox}

\section{The REPL-Driven Workflow: Interrogating the Living System}

In traditional systems, testing whether your configuration is syntactically and structurally valid requires running the full reconfigure command and hoping for the best. 

In Guix, you have \guixcmd{guix repl}.

Unlike running vanilla \texttt{guile}, \guixcmd{guix repl} initializes Guile with the complete Guix module tree, your configured channels, and the store connection automatically loaded into the environment.

\subsection{Dissecting an \texttt{operating-system} Definition Live}

Let us open \guixcmd{guix repl} and inspect our system configuration:

\begin{lstlisting}[style=guixcli]
$ guix repl
GNU Guile 3.0.9
Copyright (C) 1995-2023 Free Software Foundation, Inc.
guile> 
\end{lstlisting}

Load the necessary modules and parse your system definition file:
\begin{lstlisting}[style=guixscheme]
;; 1. Load system, service, and package modules
(use-modules (gnu)
             (gnu system)
             (gnu services)
             (ice-9 pretty-print))

;; 2. Evaluate and instantiate your configuration file
(define os (primitive-load "/etc/config.scm"))
\end{lstlisting}

The variable \texttt{os} now holds a live \texttt{<operating-system>} Scheme record object! You can now query any aspect of your machine using standard record accessors:

\begin{lstlisting}[style=guixscheme]
;; What host name is declared?
(operating-system-host-name os)
;; => "antigravity-workstation"

;; What timezone?
(operating-system-timezone os)
;; => "America/New_York"

;; List all packages declared directly in the operating-system record
(map package-name (operating-system-packages os))
;; => ("vim" "git" "htop" "tmux" "curl")
\end{lstlisting}

\subsection{Filtering and Querying Services}

The most complex part of any Guix System configuration is the service graph. In the REPL, you can inspect and query every single service running on the system:

\begin{lstlisting}[style=guixscheme]
;; Extract all services
(define svcs (operating-system-services os))

;; How many services are configured (including default base services)?
(length svcs)
;; => 47

;; Find the OpenSSH service
(use-modules (gnu services ssh))

(define ssh-matches
  (filter (lambda (s) (eq? (service-kind s) openssh-service-type))
          svcs))

(define ssh-svc (car ssh-matches))
\end{lstlisting}

Now inspect the value (the configuration record) passed to that service:
\begin{lstlisting}[style=guixscheme]
(define ssh-cfg (service-value ssh-svc))

;; Inspect individual fields
(openssh-configuration-port-number ssh-cfg)
;; => 22

(openssh-configuration-permit-root-login ssh-cfg)
;; => #f
\end{lstlisting}

\begin{protipbox}[Testing \texttt{modify-services} in the REPL]
Are you trying to write a complex \texttt{modify-services} clause to customize GDM, UDev rules, or OpenSSH, but keep getting syntax errors?

Don't test it by running \guixcmd{guix system reconfigure}! Paste your \texttt{modify-services} expression directly into \guixcmd{guix repl} and evaluate it against \texttt{svcs}. If it produces a modified list of services without throwing an exception, you know your syntax and field types are 100\% correct.
\end{protipbox}

\subsection{Inspecting Package Recipes and Inputs}

You can also use the REPL to inspect package objects, check for security updates, and evaluate dependencies:

\begin{lstlisting}[style=guixscheme]
(use-modules (gnu packages web))

;; Look up a package by specification
(define nginx (specification->package "nginx"))

;; Check package version and license
(package-version nginx)
;; => "1.24.0"

(package-license nginx)
;; => #<license: bsd-2>

;; Inspect its direct inputs
(map car (package-inputs nginx))
;; => ("openssl" "pcre" "zlib")

;; Inspect build arguments passed to the build system
(package-arguments nginx)
;; => (#:tests? #f #:configure-flags '("--with-http_ssl_module" ...))
\end{lstlisting}

\subsection{Testing G-Expressions Without Building the Universe}

In Chapter~\ref{chap:phases-and-gexp}, we learned that G-Expressions (\texttt{\#\~{}}) represent code staged for execution inside store derivations. How do you test what a G-expression expands to without running a 30-minute system build?

You can \textbf{lower} any object to a store derivation right inside the REPL:

\begin{lstlisting}[style=guixscheme]
(use-modules (guix gexp)
             (guix store)
             (guix derivations))

;; Open a live connection to the guix-daemon
(define store (open-connection))

;; Lower a package reference to its store derivation
(run-with-store store
  (lower-object (specification->package "hello")))
;; => #<derivation /gnu/store/...-hello-2.12.1.drv => ...>
\end{lstlisting}

\subsection{Guile REPL Metacommands: Your Interactive Swiss Army Knife}

Guile's REPL provides built-in metacommands (prefixed with a comma \texttt{,}) that turn the terminal into an interactive debugger:

\begin{table}[htbp]
\centering
\small
\begin{tabularx}{\textwidth}{L{0.28\textwidth} Y}
\toprule
\textbf{Metacommand} & \textbf{Function and Usage} \\
\midrule
\texttt{,describe <expr>} & View docstring, parameter list, and module location of any procedure or macro. \\
\texttt{,inspect <expr>} & Launch an interactive tree browser for any record, hash-table, vector, or closure. \\
\texttt{,trace (<call>)} & Trace procedure execution, showing arguments and return values at each stack frame. \\
\texttt{,step (<call>)} & Step interactively through expression evaluation frame-by-frame. \\
\texttt{,bt} (backtrace) & In an error break loop, print the complete call stack with local variables. \\
\texttt{,locals} & In a debugger break frame, display all bound local variables. \\
\texttt{,quit} & Exit the REPL session. \\
\bottomrule
\end{tabularx}
\caption{Essential GNU Guile REPL metacommands.}
\end{table}

For example, typing \texttt{,inspect ssh-cfg} launches an interactive curses-style inspector where you can arrow through every field of the OpenSSH configuration record in real time!

\section{The Lisp Nirvana: Emacs and Geiser Integration}

If \guixcmd{guix repl} gives you terminal-based superpowers, pairing GNU Guix with \textbf{GNU Emacs} and \textbf{Geiser} elevates your workflow to absolute nirvana.

In conventional editors (VS Code, Sublime, Vim), Scheme code is treated as static text with syntax highlighting. In Emacs with Geiser, your editor buffer is connected directly to a living Scheme evaluation process.

\begin{protipbox}[The Lisp Alchemist's Cockpit]
In an Emacs-Geiser workflow, your screen is typically split into two interactive halves:
\begin{itemize}
  \item \textbf{Left Window (Source Buffer)}: Your \texttt{config.scm} or package definition file. You write expressions, jump to definitions with \texttt{M-.}, and send definitions to the compiler.
  \item \textbf{Right Window (Guile REPL Buffer)}: The live Guix runtime. It evaluates your expressions, prints record representations, and lets you interrogate values interactively.
\end{itemize}
\end{protipbox}

\subsection{Setting Up the Environment in 60 Seconds}

You do not need to spend days configuring an elaborate Emacs setup. You can launch a complete, hermetic Emacs development environment with all required Guix extensions in a single command using \guixcmd{guix shell}:

\begin{lstlisting}[style=guixcli]
guix shell emacs \
           emacs-geiser \
           emacs-geiser-guile \
           emacs-guix \
           emacs-paredit \
           -- emacs
\end{lstlisting}

\subsection{The Geiser Workflow: Interactive Evaluation}

Open your \texttt{config.scm} or custom package file in Emacs, then start Geiser:
\begin{center}
\texttt{M-x geiser-guile}
\end{center}
This splits your window and launches a dedicated Guile REPL connected directly to your buffer.

Here are the four essential keybindings that will change the way you write Guix configurations forever:

\begin{enumerate}
  \item \textbf{Evaluate S-Expression Preceding Cursor (\texttt{C-x C-e})}:\\
    Place your cursor immediately after any closing parenthesis and press \texttt{C-x C-e}. Geiser evaluates that specific expression in the REPL and displays the result in the echo area.
  
  \item \textbf{Evaluate Entire Top-Level Definition (\texttt{C-c C-c})}:\\
    Place your cursor anywhere inside a \texttt{(define-public ...)} package or a \texttt{(define os ...)} block and press \texttt{C-c C-c}. The definition is compiled and sent to the running Guix session instantly.

  \item \textbf{Switch Between Buffer and REPL (\texttt{C-c C-z})}:\\
    Toggle instantaneously between your Scheme source buffer and the interactive REPL prompt.

  \item \textbf{Instant Documentation Lookup (\texttt{C-c C-d C-d})}:\\
    Place your cursor on any Guix procedure or macro (e.g. \texttt{modify-phases}, \texttt{service-extension}) and press \texttt{C-c C-d C-d} to display its complete documentation in a dedicated buffer.
\end{enumerate}

\subsection{The Holy Grail: Jump to Definition with \texttt{M-.}}

This is the single feature that transforms Guix from an overwhelming wall of Scheme code into an intuitive, self-documenting playground:

Place your cursor on \textbf{any} symbol in your buffer:
\begin{itemize}
  \item An unfamiliar service type: \texttt{dhcp-client-service-type}
  \item A configuration record: \texttt{tor-configuration}
  \item A build system helper: \texttt{substitute*}
  \item A package definition: \texttt{git}
\end{itemize}

Press:
\begin{center}
\Large\textbf{\texttt{M-.}} \normalsize\quad (\texttt{geiser-edit-symbol-at-point})
\end{center}

Emacs will \textbf{instantly jump to the exact source file and line number in the GNU Guix codebase} where that symbol was defined! 

You can read how the service is constructed, see what default arguments it expects, inspect its Shepherd action scripts, and copy patterns directly from the core maintainers.

When you are finished reading, press:
\begin{center}
\Large\textbf{\texttt{M-,}} \normalsize\quad (\texttt{geiser-pop-symbol-stack})
\end{center}
and Emacs jumps right back to your personal configuration file. You will never need to hunt through external wiki pages or outdated blog posts again.

\subsection{Browsing Guix Visually with \texttt{guix.el}}

The \texttt{emacs-guix} package provides a rich visual interface to Guix package management directly inside Emacs buffers:

\begin{itemize}
  \item \texttt{M-x guix-packages-by-name}: Search for packages with an interactive tabular browser showing descriptions, homepages, licenses, and outputs.
  \item \texttt{M-x guix-generations}: View your system and profile generation history. You can compare and diff any two generations with a single keystroke (\texttt{d}) to see exactly which packages were added, removed, or upgraded!
  \item \texttt{M-x guix-system-generations}: Inspect, activate, or delete Guix System generations visually.
\end{itemize}

\section{Runtime Triage: Shepherd Services and Disposable Virtual Machines}

What happens when your package builds, your \texttt{config.scm} evaluates cleanly, \guixcmd{guix system reconfigure} succeeds, but your computer reboots and the web server or desktop manager fails to start?

\subsection{Diagnosing Shepherd Service Failures}

GNU Shepherd (Chapter~\ref{chap:shepherd}) is your PID 1 supervisor. When a service fails to launch, query it with \guixcmd{herd}:

\begin{lstlisting}[style=guixcli]
# Check status of all running and stopped services
herd status

# Get deep diagnostics on a failing service
herd detailed-status nginx
\end{lstlisting}

\guixcmd{herd detailed-status} displays:
\begin{itemize}
  \item Whether the service is currently running or stopped.
  \item The PID of the running daemon.
  \item Recent exit statuses and timestamps (e.g. \texttt{exited with code 1 at 14:02:11}).
  \item The \textbf{Respawn count}: whether Shepherd tried to restart the service and gave up because it was crashing in an infinite loop.
\end{itemize}

\subsection{Inspecting Shepherd and System Logs}

Shepherd records all service stdout, stderr, and lifecycle transitions in \texttt{/var/log/}:

\begin{lstlisting}[style=guixcli]
# View Shepherd daemon events
sudo less /var/log/shepherd.log

# View system messages and kernel syslog
sudo less /var/log/messages

# Follow logs in real time while starting a service
sudo tail -f /var/log/shepherd.log
\end{lstlisting}

\subsection{Zero-Risk Testing with Ephemeral Virtual Machines}

Never test a radical system reconfiguration (such as modifying your display manager, core PAM authentication rules, or disk encryption) directly on bare metal if you cannot afford downtime.

Guix System has a miraculous built-in superpower: \guixcmd{guix system vm}.

\begin{lstlisting}[style=guixcli]
# Compile your configuration into a self-contained QEMU virtual machine script
guix system vm /etc/config.scm
\end{lstlisting}

This command evaluates your configuration, builds any missing packages, and produces a single executable script in \texttt{/gnu/store/...-run-vm.sh}:

\begin{lstlisting}[style=guixcli]
# Launch your proposed operating system in a window!
$(guix system vm /etc/config.scm) -m 2048 -smp 2 -net nic -net user,hostfwd=tcp::2222-:22
\end{lstlisting}

\begin{alchemybox}[The Power of the Ephemeral VM]
When you boot this virtual machine:
\begin{itemize}
  \item It boots your \textbf{exact} declared kernel, Shepherd services, user accounts, and desktop environment inside QEMU.
  \item It shares the immutable host store (\texttt{/gnu/store}) read-only via 9p virtio filesystem sharing, so starting the VM takes less than 5 seconds!
  \item Any disk writes in the VM are redirected to an ephemeral QEMU overlay file. Your host machine is 100\% protected.
  \item You can SSH into your VM on port 2222 (\texttt{ssh -p 2222 user@localhost}) to test daemons, verify network configurations, and validate Shepherd scripts before running \guixcmd{guix system reconfigure} on bare metal.
\end{itemize}
\end{alchemybox}

\section{Real-World Case Studies: Finding Your Way in the Fog}

Abstract theory is comforting, but true diagnostic wisdom is forged in the fires of actual failure. Here are six battle-tested case studies representing the most frequent stumbling blocks encountered by newcomers, along with the exact step-by-step investigative workflows to resolve them.

\subsection{Case Study 1: The Missing Module and the Unbound Variable Mystery}

\textbf{The Setup}: You are configuring a developer workstation. You decide to add the Bluetooth service and a Tor daemon to your \texttt{services} list in \texttt{/etc/config.scm}:

\begin{lstlisting}[style=guixscheme]
(operating-system
  ;; ...
  (services
    (cons* (service bluetooth-service-type)
           (service tor-service-type)
           %desktop-services)))
\end{lstlisting}

\textbf{The Error}: You run \guixcmd{sudo guix system reconfigure /etc/config.scm} and get slapped with:
\begin{lstlisting}[style=guixcli]
/etc/config.scm:42:19: error: bluetooth-service-type: unbound variable
hint: Did you forget `(use-modules (gnu services desktop))'?
\end{lstlisting}
You dutifully add \texttt{(use-modules (gnu services desktop))} at the top of the file and re-run. Now Guix complains:
\begin{lstlisting}[style=guixcli]
/etc/config.scm:43:19: error: tor-service-type: unbound variable
\end{lstlisting}
Notice that for \texttt{tor-service-type}, \textbf{Guix offers no hint at all!}

\textbf{The Investigation}: Where on earth is \texttt{tor-service-type} defined? Do you comb through 1,200 source files in Guix?
No! You have three immediate tools:

\begin{enumerate}
  \item \textbf{The Geiser Apropos Search (in Emacs)}:\\
    Press \texttt{C-c C-d C-a} (\texttt{geiser-apropos}) and type \texttt{tor-service}. Within half a second, Emacs scans all loaded and discovered modules and displays:
    \begin{lstlisting}[style=guixscheme]
    tor-service-type: <service-type> in (gnu services networking)
    \end{lstlisting}

  \item \textbf{The Fast Terminal \texttt{grep} over Guix Sources}:\\
    If you are in a terminal, query the Scheme load path for services:
    \begin{lstlisting}[style=guixcli]
    # Locate the services directory inside the store
    guix repl -e '(display (search-path %load-path "gnu/services.scm"))'
    \end{lstlisting}
    Or simply grep the installed module tree:
    \begin{lstlisting}[style=guixcli]
    find $(guix repl -e '(display (dirname (search-path %load-path "gnu/services.scm")))') \
      -name "*.scm" | xargs grep -l "define-public tor-service-type"
    \end{lstlisting}
    It reports: \texttt{gnu/services/networking.scm}!
\end{enumerate}

\textbf{The Resolution}: Add \texttt{(gnu services networking)} to your \texttt{use-modules} declaration:
\begin{lstlisting}[style=guixscheme]
(use-modules (gnu)
             (gnu services desktop)     ; for bluetooth-service-type
             (gnu services networking))  ; for tor-service-type
\end{lstlisting}

\subsection{Case Study 2: Record Archeology---Taming Cryptic Field Errors}

\textbf{The Setup}: You want to run the OpenSSH daemon on a custom port and disable password authentication. You write:

\begin{lstlisting}[style=guixscheme]
(service openssh-service-type
  (openssh-configuration
    (port 2222)
    (password-authentication? "no")))
\end{lstlisting}

\textbf{The Error}: You reconfigure, and Guix crashes with:
\begin{lstlisting}[style=guixcli]
ice-9/boot-9.scm:1685:16: In procedure struct-vtable: Wrong type argument: #f
error: invalid field specifier `port`
\end{lstlisting}
Or, if you fix the field name to something plausible, you get:
\begin{lstlisting}[style=guixcli]
error: value "no" does not match predicate boolean?
\end{lstlisting}

\textbf{The Investigation}: What are the exact field names and types for \texttt{openssh-configuration}?
Instead of guessing or reading third-party blog posts from 2019, inspect the record definition directly:

\begin{enumerate}
  \item \textbf{In Emacs}: Put cursor on \texttt{openssh-configuration} and press \texttt{M-.}. You land directly on line 140 of \texttt{gnu/services/ssh.scm}.
  \item \textbf{In Terminal}: Launch \guixcmd{guix repl} and inspect the record type:
    \begin{lstlisting}[style=guixscheme]
    (use-modules (gnu services ssh))
    ,describe openssh-configuration
    \end{lstlisting}
    Even better, instantiate the default configuration and inspect it:
    \begin{lstlisting}[style=guixscheme]
    ,inspect (openssh-configuration)
    \end{lstlisting}
    Guile opens an interactive record browser displaying every field, its name, and its default value:
    \begin{lstlisting}[style=guixcli]
    [0] port-number: 22
    [1] permit-root-login: #f
    [2] password-authentication?: #t
    [3] public-key-authentication?: #t
    [4] x11-forwarding?: #f
    \end{lstlisting}
\end{enumerate}

\textbf{The Revelation}:
\begin{itemize}
  \item The field is not \texttt{port}; it is \texttt{port-number}, expecting an exact integer.
  \item The field is \texttt{password-authentication?}, expecting a Scheme boolean (\texttt{\#t} or \texttt{\#f}), not a string like \texttt{"no"}.
\end{itemize}

\textbf{The Resolution}:
\begin{lstlisting}[style=guixscheme]
(service openssh-service-type
  (openssh-configuration
    (port-number 2222)
    (password-authentication? #f)))
\end{lstlisting}

\subsection{Case Study 3: The Service Collision: ``Service X provided more than once''}

\textbf{The Setup}: You just bought a graphic drawing tablet or hardware security key (e.g. YubiKey). To make it work, you must add custom UDev rules. You look up \texttt{udev-service-type} and write:

\begin{lstlisting}[style=guixscheme]
(operating-system
  ;; ...
  (services
    (cons* (service udev-service-type
             (udev-configuration
               (rules (list my-custom-udev-rules))))
           %desktop-services)))
\end{lstlisting}

\textbf{The Error}:
\begin{lstlisting}[style=guixcli]
guix system: error: service 'udev' provided more than once
\end{lstlisting}

\begin{tearsbox}[The Mental Model Collision]
Why did this happen? In imperative distros, you just drop a file into \texttt{/etc/udev/rules.d/}. In Guix, \texttt{\%desktop-services} already instantiates and configures a system-wide UDev daemon. 

When you declared \texttt{(service udev-service-type ...)}, you told Guix to launch a \textbf{second, competing UDev daemon} managing the same hardware! Guix's service graph validator detects this collision and aborts before you render your system unbootable.
\end{tearsbox}

\textbf{The Investigation in the REPL}: Let us verify where the existing UDev service lives:
\begin{lstlisting}[style=guixscheme]
(use-modules (gnu) (gnu services) (gnu services base))
(define svcs %desktop-services)

;; Find all services that are of kind udev-service-type
(filter (lambda (s) (eq? (service-kind s) udev-service-type)) svcs)
;; => (#<service udev 7f3b890a2c00>)
\end{lstlisting}
Indeed, UDev is already a core member of \texttt{\%desktop-services}!

\textbf{The Resolution}: In Guix, you do not add a duplicate service; you **extend** the existing one. Guix provides a dedicated helper for this exact scenario:

\begin{lstlisting}[style=guixscheme]
(use-modules (gnu services base))

(operating-system
  ;; ...
  (services
    (cons* (udev-rules-service 'my-drawing-tablet my-custom-udev-rules)
           %desktop-services)))
\end{lstlisting}
The \texttt{udev-rules-service} helper emits a service extension that automatically injects your rules into the existing UDev daemon without collision!

\subsection{Case Study 4: The Build Crime Scene: Missing Headers and the Phantom \texttt{pkg-config}}

\textbf{The Setup}: You are packaging a C utility called \texttt{fast-calc} that depends on SQLite. You write your package definition:

\begin{lstlisting}[style=guixscheme]
(define-public fast-calc
  (package
    (name "fast-calc")
    (version "1.0")
    (source (origin ...))
    (build-system gnu-build-system)
    (inputs (list sqlite))
    (synopsis "Ultra fast calculator")
    ...))
\end{lstlisting}

\textbf{The Error}: You run \guixcmd{guix build -f fast-calc.scm} and get:
\begin{lstlisting}[style=guixcli]
starting phase `configure'
checking for gcc... gcc
checking for sqlite3... no
configure: error: Package requirements (sqlite3 >= 3.0) were not met:
No package 'sqlite3' found
Consider adjusting the PKG_CONFIG_PATH environment variable...
error: in phase 'configure': returned #f
\end{lstlisting}

\textbf{The Panic}: You scream: \emph{``I literally put \texttt{sqlite} in \texttt{inputs}! Why does configure say sqlite3 is not found?!''}

\textbf{The Forensic Investigation}:
Run the build with \guixcmd{-K} to capture the sandbox:
\begin{lstlisting}[style=guixcli]
guix build -K -f fast-calc.scm
\end{lstlisting}
Guix keeps the directory in \texttt{/tmp/guix-build-fast-calc-1.0.drv-0}.
Now, step inside the exact environment:
\begin{lstlisting}[style=guixcli]
guix shell -D -f fast-calc.scm -- bash
cd /tmp/guix-build-fast-calc-1.0.drv-0
source ./environment-variables
\end{lstlisting}

Now interrogate the environment:
\begin{lstlisting}[style=guixcli]
# 1. Does SQLite exist in the store?
echo $CPATH
# => /gnu/store/...-sqlite-3.40.1/include

# 2. Does the .pc file exist in PKG_CONFIG_PATH?
echo $PKG_CONFIG_PATH
# => /gnu/store/...-sqlite-3.40.1/lib/pkgconfig

# 3. Can pkg-config find it?
pkg-config --modversion sqlite3
# => bash: pkg-config: command not found!
\end{lstlisting}

\textbf{The Eureka Moment}: SQLite was present, its headers were present, and its \texttt{.pc} file was in \texttt{PKG\_CONFIG\_PATH}. But the \texttt{configure} script used the \texttt{pkg.m4} macro to look for SQLite, and \textbf{\texttt{pkg-config} itself was never installed in the build container!}

Because \texttt{pkg-config} is a build tool executed on the build host (not a library linked into the target binary), it belongs in \texttt{native-inputs}.

\textbf{The Resolution}:
\begin{lstlisting}[style=guixscheme]
(define-public fast-calc
  (package
    (name "fast-calc")
    ;; ...
    (native-inputs (list pkg-config))  ; <--- The missing link!
    (inputs (list sqlite))
    ;; ...))
\end{lstlisting}
Re-running \guixcmd{guix build -f fast-calc.scm} now succeeds in 4 seconds flat.

\subsection{Case Study 5: Post-Reconfigure Triage: The Crashing Shepherd Daemon}

\textbf{The Setup}: You configured an Nginx web server on your Guix machine to host an internal dashboard:

\begin{lstlisting}[style=guixscheme]
(service nginx-service-type
  (nginx-configuration
    (server-blocks
      (list (nginx-server-configuration
              (listen '("80"))
              (server-name '("localhost"))
              (root "/srv/http"))))))
\end{lstlisting}

\textbf{The Symptom}: \guixcmd{sudo guix system reconfigure /etc/config.scm} completes cleanly. No build errors, no Scheme errors. But when you run \texttt{curl http://localhost}, you receive:
\begin{lstlisting}[style=guixcli]
curl: (7) Failed to connect to localhost port 80: Connection refused
\end{lstlisting}

\textbf{The Investigation}:
Check Shepherd service status:
\begin{lstlisting}[style=guixcli]
sudo herd status nginx
\end{lstlisting}
Shepherd reports:
\begin{lstlisting}[style=guixcli]
Status of nginx:
  It is stopped.
  It is enabled.
  Provides: (nginx httpd)
  Requires: (user-processes loopback)
\end{lstlisting}
Try starting it manually:
\begin{lstlisting}[style=guixcli]
sudo herd start nginx
\end{lstlisting}
Shepherd prints:
\begin{lstlisting}[style=guixcli]
herd: error: Service nginx could not be started.
\end{lstlisting}

Why? Ask Shepherd for detailed forensic telemetry:
\begin{lstlisting}[style=guixcli]
sudo herd detailed-status nginx
\end{lstlisting}
Output:
\begin{lstlisting}[style=guixcli]
  Last exit code: 1
  Respawn limit hit: 5 crashes in 3 seconds. Disabled respawning.
\end{lstlisting}

Now check the Shepherd daemon log and Nginx error log:
\begin{lstlisting}[style=guixcli]
sudo tail -n 25 /var/log/shepherd.log
sudo cat /var/log/nginx/error.log
\end{lstlisting}
The error log reveals the truth in plain English:
\begin{lstlisting}[style=guixcli]
2026/11/14 10:15:22 [emerg] 1421#1421: bind() to 0.0.0.0:80 failed (98: Address already in use)
\end{lstlisting}

\textbf{The Resolution}:
Run \guixcmd{sudo ss -tulpn | grep :80} to see who owns port 80. You find a stray Apache or lighttpd container running in the background! Terminate the rogue process and run \guixcmd{sudo herd start nginx}. The service starts instantly with \texttt{It is started [PID 1842]}.

\subsection{Case Study 6: The Accidental 4-Hour Compile: Diagnosing Substitute Cache Misses}

\textbf{The Setup}: You want to install a popular desktop application:
\begin{lstlisting}[style=guixcli]
guix install inkscape
\end{lstlisting}
Instead of downloading pre-built binaries in 30 seconds, your laptop fans spin up to deafening speed. Your terminal prints:
\begin{lstlisting}[style=guixcli]
The following derivation will be built:
  /gnu/store/...-gcc-11.3.0.drv
building /gnu/store/...-gcc-11.3.0.drv...
[  1%] Building CXX object ...
\end{lstlisting}
Your machine is compiling GCC and LLVM from source!

\begin{footgunbox}[The Gentoo Nightmare Scenario]
Newcomers often panic here: \emph{``I thought Guix had binary substitutes! Why is it compiling the entire C compiler toolchain on my dual-core laptop?!''}
\end{footgunbox}

\textbf{The Diagnostic Workflow}:
Stop the build immediately with \texttt{Ctrl-C}. Then investigate why substitutes were not found:

\begin{enumerate}
  \item \textbf{Always Dry-Run First (\texttt{-n})}:\\
    Before running any install, build, or reconfigure, pass \guixcmd{-n} (or \guixcmd{--dry-run}):
    \begin{lstlisting}[style=guixcli]
    guix install inkscape -n
    \end{lstlisting}
    Guix will tell you upfront:
    \begin{itemize}
      \item \texttt{124.5 MB would be downloaded}: Pre-built binaries exist. Safe to proceed!
      \item \texttt{The following derivations will be built}: No substitute exists. You will be compiling from source!
    \end{itemize}

  \item \textbf{Check Substitute Server Health with \texttt{guix weather}}:\\
    \begin{lstlisting}[style=guixcli]
    guix weather inkscape
    \end{lstlisting}
    This queries the official build farm and reports whether binary substitutes are currently available for your architecture.

  \item \textbf{Identify Substitute URL Configuration}:\\
    Check whether your \texttt{guix-daemon} is configured with authorized substitute keys and servers:
    \begin{lstlisting}[style=guixcli]
    # Check authorized keys
    ls -l /etc/guix/acl
    \end{lstlisting}
    If \texttt{/etc/guix/acl} is missing or empty, the daemon refuses to download remote substitutes for security reasons! Authorize the official build farm:
    \begin{lstlisting}[style=guixcli]
    PUBKEY=$(guix repl -e \
      '(display (search-path %load-path "gnu/packages/substitutes/bordeaux.guix.gnu.org.pub"))')
    sudo guix archive --authorize < "$PUBKEY"
    \end{lstlisting}

  \item \textbf{Diagnose Channel Drift}:\\
    If you ran \guixcmd{guix pull} five minutes after a maintainer committed a change to the core toolchain, the continuous integration servers might still be compiling the binaries in the queue.
    
    Using \guixcmd{guix describe}, check your current commit. If you pulled too close to the bleeding edge, use \guixcmd{guix time-machine} to pin to a commit from 24 hours earlier where 100\% of binary substitutes are cached.
\end{enumerate}

\section{The Diagnostic Alchemist's Field Guide}

When things go bump in the night, consult these diagnostic reference matrices:

\subsection*{Evaluation and Build Diagnostics}

\begin{table}[htbp]
\centering
\small
\begin{tabularx}{\textwidth}{L{0.27\textwidth} L{0.33\textwidth} Y}
\toprule
\textbf{Symptom} & \textbf{Diagnostic Tool} & \textbf{Alchemical Cure} \\
\midrule
Cryptic Guile backtrace mentioning \texttt{boot-9.scm} &
  \guixcmd{guix system build -d -v 2} \newline \guixcmd{guix repl} &
  Identify evaluation-time syntax or record type errors; inspect stack frames from bottom to top. \\
\addlinespace
Build fails during \texttt{make} or \texttt{make check} in chroot &
  \guixcmd{guix build -K <pkg>} \newline \guixcmd{guix shell -D <pkg>} &
  Inspect \texttt{/tmp/guix-build-...}; source \texttt{environment-variables} and debug interactively. \\
\addlinespace
Missing C library during \texttt{configure} despite \texttt{inputs} &
  \guixcmd{guix build -K} \newline \texttt{environment-variables} &
  Check if \texttt{pkg-config} or build tools are missing from \texttt{native-inputs}. \\
\addlinespace
Unsure why package pulls in huge unwanted dependency &
  \guixcmd{guix size <pkg>} \newline \guixcmd{guix graph --path} &
  Trace shortest dependency path; eliminate accidental propagated inputs or hardcoded store references. \\
\addlinespace
Cannot find where a package or service is defined &
  \guixcmd{guix edit <pkg>} \newline \texttt{M-.} in Emacs / Geiser &
  Jump directly to definition line in the Guix source tree. \\
\bottomrule
\end{tabularx}
\caption{Evaluation and Build Diagnostics.}
\end{table}

\subsection*{Service, Runtime, and Substitute Diagnostics}

\begin{table}[htbp]
\centering
\small
\begin{tabularx}{\textwidth}{L{0.27\textwidth} L{0.33\textwidth} Y}
\toprule
\textbf{Symptom} & \textbf{Diagnostic Tool} & \textbf{Alchemical Cure} \\
\midrule
Syntax error in \texttt{modify-services} or system record &
  \guixcmd{guix repl} \newline \texttt{primitive-load} &
  Load \texttt{config.scm} in the REPL; evaluate record accessors and service filters interactively. \\
\addlinespace
``Service X provided more than once'' conflict &
  \guixcmd{guix repl} \newline \texttt{operating-system-services} &
  Avoid duplicate instantiation; extend the existing service via helpers (e.g. \texttt{udev-rules-service}). \\
\addlinespace
System service fails to start or crashes on boot &
  \guixcmd{herd detailed-status <svc>} \newline \texttt{/var/log/shepherd.log} &
  Check service respawn counts, exit codes, port collisions, and daemon logs. \\
\addlinespace
Unexpected 4-hour source compilation of core toolchain &
  \guixcmd{guix weather} \newline \guixcmd{guix build -n} &
  Check substitute server ACL keys; avoid un-cached bleeding-edge channel commits with \guixcmd{time-machine}. \\
\addlinespace
Fear of breaking workstation on reconfigure &
  \guixcmd{guix system vm config.scm} &
  Boot proposed system in a 5-second QEMU sandbox before touching bare metal. \\
\bottomrule
\end{tabularx}
\caption{Service, Runtime, and Substitute Diagnostics.}
\end{table}

With these tools and case studies at your command, the parentheses are no longer an impenetrable barrier. They are a transparent window into the living soul of your operating system.
